\documentclass[10pt,a4paper]{article}
\usepackage[margin=2cm,headheight=14pt]{geometry}
\usepackage{xcolor}
\usepackage{tcolorbox}
\tcbuselibrary{skins,breakable}
\usepackage{enumitem}
\usepackage{fancyhdr}
\usepackage{siunitx}
\DeclareSIUnit\molar{M}
\usepackage{titlesec}
\usepackage{microtype}
\usepackage[hidelinks]{hyperref}

\definecolor{ink}{RGB}{35,60,95}
\definecolor{paper}{RGB}{252,250,243}
\setlength{\parindent}{0pt}

\pagestyle{fancy}
\fancyhf{}
\lhead{\small\textsf{Lab notebook: Project Name}}
\rhead{\small\textsf{Researcher Name}}
\cfoot{\small\thepage}

% Each dated entry starts with \entry{date}{short title}.
\titleformat{\section}{\normalfont\sffamily\large\bfseries\color{ink}}{}{0pt}{}
\newcommand{\entry}[2]{\clearpage\section{#1 \textbar\ #2}}

% Boxes for the recurring parts of an entry.
\newtcolorbox{aim}{enhanced,breakable,colback=paper,colframe=ink,boxrule=0.6pt,arc=1mm,title=Aim,fonttitle=\sffamily\bfseries}
\newtcolorbox{observation}{enhanced,breakable,colback=white,colframe=ink!50,boxrule=0.4pt,arc=1mm,title=Observations,fonttitle=\sffamily\bfseries,coltitle=black,colbacktitle=ink!10}
\newtcolorbox{nextsteps}{enhanced,breakable,colback=ink!5,colframe=ink!5,arc=1mm,left=2mm}

\begin{document}

\begin{center}
{\sffamily\Huge\bfseries\color{ink} Research Notebook\par}
\vspace{0.5em}
{\large Project: Enzyme kinetics of variant lactase\par}
\vspace{0.3em}
Owner: Researcher Name \quad Lab: Room 2.14 \quad Started: 5 January 2026
\end{center}

\vspace{1em}
\begin{tcolorbox}[colback=paper,colframe=ink,title=How to use this notebook,fonttitle=\sffamily\bfseries]
Add a new entry for every working day. Record what you planned, what you actually did, raw observations and what to try next. Never delete a failed result: note why it failed.
\end{tcolorbox}

\tableofcontents

\entry{2026-01-05}{Buffer preparation}

\begin{aim}
Prepare phosphate buffer for the week's assays and check its pH.
\end{aim}

\textbf{Materials}
\begin{itemize}[nosep]
  \item \SI{0.1}{\molar} sodium phosphate, \SI{500}{\milli\litre}
  \item pH meter, calibrated at 09:10
\end{itemize}

\textbf{Procedure}
\begin{enumerate}[nosep]
  \item Dissolved salts in \SI{400}{\milli\litre} deionised water.
  \item Adjusted to pH 7.0 and topped up to volume.
\end{enumerate}

\begin{observation}
Final pH 7.02 at \SI{22}{\celsius}. Solution clear. Stored at \SI{4}{\celsius}, bottle labelled B-01.
\end{observation}

\begin{nextsteps}
\textbf{Next:} run the substrate dilution series tomorrow. Order more ONPG.
\end{nextsteps}

\entry{2026-01-06}{Substrate dilution series}

\begin{aim}
Measure initial reaction rates at six substrate concentrations.
\end{aim}

\begin{center}
\begin{tabular}{ccc}
\hline
Tube & [S] (\si{\milli\molar}) & $A_{420}$ after \SI{60}{\second} \\
\hline
1 & 0.5 & 0.081 \\
2 & 1.0 & 0.142 \\
3 & 2.0 & 0.231 \\
4 & 4.0 & 0.330 \\
\hline
\end{tabular}
\end{center}

\begin{observation}
Tube 4 showed bubbles in the cuvette; repeat before fitting. Sketch or paste plots below.
\end{observation}

\fbox{\parbox[c][4cm][c]{\dimexpr\linewidth-2\fboxsep-2\fboxrule}{\centering\textsf{Figure placeholder: rate against [S]}}}

\begin{nextsteps}
\textbf{Next:} repeat tube 4, then fit Michaelis and Menten parameters.
\end{nextsteps}

\entry{YYYY-MM-DD}{Entry title}
Copy this empty entry for each new day.

\begin{aim}
What do you plan to do today?
\end{aim}

\begin{observation}
What happened? Include numbers, units and anything unexpected.
\end{observation}

\begin{nextsteps}
\textbf{Next:}
\end{nextsteps}

\end{document}
